\documentclass[10pt]{article}

\usepackage[utf8]{inputenc}

\usepackage[T1]{fontenc}

\usepackage{amsmath}

\usepackage{amsfonts}

\usepackage{amssymb}

\usepackage[version=4]{mhchem}

\usepackage{stmaryrd}

\usepackage{mathrsfs}

\usepackage{graphicx}

\usepackage[export]{adjustbox}

\graphicspath{ {./images/} }



\begin{document}

$E \otimes T_{X}$, заданная формулой



\[

H_{x}(v \otimes u)=h_{x}\left(v, \Omega_{x}(u) v\right),

\]



где $x \in X, v \in E_{x}, u \in T_{X}, x$, положительно определена. Примеры: касательное расслоение $T_{P_{n}}$ к проективному пространству $P^{n}$ положительно, но при $n>1$ не является положительным в смысле Накано; расслоение на комплексные прямые над $P^{n}$, определяемое гиперплоскостью, положительно.



Любое факторрасслоение положительного векторного расслоения положительно. Если $E^{\prime}, E^{\prime \prime}-$ положительные (положительные в смысле Накано) расслоения, то $E^{\prime} \oplus E^{\prime \prime}$ и $E^{\prime} \otimes E^{\prime \prime}$ положительны (положительны в смысле Накано).



Понятие «П. р.» было введено в связи с Кодаиры теоремой об обращении в нуль для случая расслоений на комплексные прямые, а затем обобщено на произвольные расслоения. Несколько позже, в связи с вопросом о существовании вложения в проективное пространство, были выделены понятия слабо положительного и слабо отрицательного расслоений.



Голоморфное векторное расслоение $E$ над компактным комплексным пространством $X$ наз. с л а б о о т р и ц а т е ль ны м, если его нулевое сечение обладает строго псевдовыпуклой окрестностью в $E$, т. е. является исключительным аналитич. множеством. Расслоение $E$ наз. с л а 6 о по л о ж и т е л ь н ы м, если сопряженвое расслоение $E^{*}$ слабо отрицательно. В случае, когда $X-$ риманова поверхность, понятия слабо положительного и П. р. совпадают [5]. В общем случае из положительности вытекает слабая положительность; примеров слабо положительных, но не положительных расслоений пока (1983) не известно.



Слабая положительность расслоения $E \rightarrow X$ равносильна каждому из следующих свойств: для любого когерентного аналитич. пучка $\mathscr{F}$ на $X$ существует такое $m_{0}>0$, что пучок $\mathscr{F} \otimes S^{m} \mathscr{O}$ при $m \geqslant m_{0}$ порождается глобальными сечениями; для любого когерентного аналитич. пучка $\mathscr{F}$ на $X$ существует такое $m \geqslant 0$, पто



\[

H^{q}\left(X, \mathscr{F} \otimes S^{m} \mathscr{O}\right)=0

\]



для всех $q \geqslant 1$ (см. [3], [4]). Через ஜீ здесь обозначается пучок ростков голоморфных сечений расслоения $E$. Слабо положительные расслоения аналогичны, таким образом, обильнылм векторныл расслоениям из алгебраич. геометрии и иногда наз. о б и л ь н ы м и ан алит и ч еским и расслоения ми. Слабо положительное расслоение над пространством $X$ естественным образом определяет вложение пространства $X$ в многообразие Грассмана и тем самым в проективное пространство.



Понятия положительного, отрицательного, слабо положительного и слабо отрицательного расслоений естественным образом обобщаются также на случай линейных пространств над комплексным пространством X (см. Векторное аналитическое расслоение).



См. также Oтрицательное расслоение.



Лит : [1] У ә л л с Р., Дифференциальное исчисление на комплексных многообразиях, пер. с авгл., М., 1976; 12] Ч ж ә н ь Ші ә н-ш ә н ь, Комплексные многообразия, пер. с англ., М., 1961; [3] 'Итоги науки и техники. Алгебра. Топология. Геомerpur, T. 15, M., 1977, c. 93-171; [4] Sc hn e id e r M., "70; [5] U m e m u r a H., "Nagoya Math. J.", 1973, v. 52, 97-128. A. Л. Oнuщuк.



ПОЛОЖИТЕЛЬНЫЙ КОНУС - подмножество $K$ действительного векторного пространства $E$, удовлетворяющее следующим условиям:



\begin{enumerate}

  \item если $x, y \in K, \alpha, \beta \geqslant 0$, то $\alpha x+\beta y \in K$;



  \item $K \cap(-K)=0$.



\end{enumerate}



П. к. определяет полуупорядочевие в $E$ : по определению, $x \prec y$, если $y-x \in K$. Пусть $E$ - банахово пространство. Если $E-$ замкнутый в о с п р о и 3 в о д я щ й П. К. (т. е. такой П. к., каждый вектор $z \in E$ к-рого представим в виде $z=x-y$, где $x, y \in K)$, то $\|x\|+\|y\| \leqslant M\|z\|, M$ не зависит от $z$. Т е л е с ны й, т. е. имеющий внутренние точки, П. к. является воспроизводящим.



Пусть $E^{*}$ - сопряженное пространство к банахову пространству $E$. Если $K \subset E$ - воспроизводящий П. к., то множество $K^{*} \subset E^{*}$ положительных (относительно П. к.) функционалов (т. е. таких $f$, что $f(x) \geqslant 0$ при $x \in K$ ) является П. к. (это т. н. с о п р я ж е в ны й к о н у с). П. к. $K$ восстанавливается на $K^{*}$, а именно:



\[

K=\left\{x \in E \mid f(x) \geqslant 0 \text { при } f \in K^{*}\right\} .

\]



Если $K$ - телесный П. к., то его внутренность совпадает с



\[

\left\{x \in E \mid f(x)>0 \text { при } f \in K^{*}, f \neq 0\right\} .

\]



Конус в банаховом пространстве $E$ наз. н о р м а л ьны м, если найдется такое $\delta>0$, что $\|x+y\| \geqslant \delta(\|u\|+$ $+\|y\|)$ при $x, y \in K$. П. к. нормален тогда и только тогда, когда сопряженный конус $K^{*}$ является воспроизводящим. Если $K-$ воспроизводяпий конус, то сопряженный конус $K^{*}$ нормален.



Конус $K$ наз. м и н и ә д р а л ь н ы м, если каждая пара элементов $x, y \in E$ имеет точную верхнюю грань $z=\sup (x, y)$ (то есть $z \geqslant x, y$ и для любого $z_{1} \in E$ из $z_{1} \geqslant x, y$ вытекает $z_{1} \geqslant z$ ). Если П. к. правилен и миниэдрален, то всякое счетное ограниченное подмножество имеет точную верхнюю грань.



Лum.: [1] К р а с н о с е Ј в с к и й М. А., Положительные решения операторных уравнений, М., 1962. В. Ломоносов.



\includegraphics[max width=\textwidth, center]{2023_07_08_71c1b403e11941f4f02cg-1}

тельное отоб т а жение,-1) П. о. в гиль6 е р то в м п р о с т р а н с т е - линейный оператор $A$, для к-рого соответствующая квадратичная форма $(A x, x)$ неотрицательна. П. о. необходимо симметричен и допускает самосопряженное расширение, также являющееся П. о. Самосопряженный оператор $A$ является П. о. тогда и только тогда, когда выполняется любое из следующих условий: а) $A=B^{*} B$, где $B$ - замкнутый оператор; б) $A=B^{2}$, где $B$ - самосопряженный оператор; в) спектр $A$ содержится в $(0, \infty)$. Совокупность ограниченных П. о. в гильбертовом пространстве образует конус в алгебре всех ограниченных операторов.



\begin{enumerate}

  \setcounter{enumi}{1}

  \item П. о. в п р о с т р а н с т в е с кон у с о мотображение векторного пространства $X$ в себя, сохраняющее выделенный в $X$ конус $K$. Интегральные операторы с положительными ядрами в различных функциональных пространствах с выделенными конусами положительных функций являются линейными П. о. При нек-рых дополнительных условиях на геометрию конуса $K$ и действие П. о. $A$ удается установить существование $A$-собственных векторов из $X$ (соответствующие собственные значения наз. П о 3 ит и в в ы м и, или в е д ущ и м и, они превосходят абсолютные величины всех остальных собственных значений). Напр., доказано (см. [3]), что если $A$ вполне непрерывный П. о. с ненулевым спектром, то его спектральный радиус является позитивным собетвенным значением. Условие компактности можно заменить условиями на поведение резольвенты (см. [4]).

\end{enumerate}



В случае велинейных П. о. исследуется вопрос о существовании неподвижной точки оператора (т. е. решения уравнения $A x=x$ ) и о возможности нахождения этой точки как предела определенных рекуррентных последователей.



Нек-рые результаты теории П. о. могут быть перенесены на операторы, оставляющие инвариантными выделенные подмножества более общей ігрироды, чем конусы (см. [5]).





\end{document}